% =====================================================================
% Discussion, future work, conclusion  (~280 words)
% =====================================================================
\section{Discussion and Future Work}
\label{sec:futurework}

The principal advanced-feature claims that drive this submission
above the 60\,\% core-only ceiling are: (i)~a custom UART framing
protocol with start-byte re-sync, XOR checksum and link-loss
watchdog, hosted above the HC-05 SPP tunnel; (ii)~ISR-driven 100\,Hz
PID control on every joint with anti-windup; (iii)~an explicit FSM
for the autonomous mode with a \texttt{SAFE} fault state;
(iv)~a multi-MCU heterogeneous architecture in which an ESP32 (later
a Raspberry Pi) handles computer vision while the Mega retains all
real-time motor and feedback responsibilities; and (v)~quantitative
bench instrumentation that drives iterative optimisation rather than
ad-hoc demonstration.

Several extensions remain. The Mega super-loop is structured to
allow a deferred move to ArduinoFreeRTOS, which would convert the
joystick sampling, packet RX, motor control and feedback PID into
explicit tasks with deterministic priorities; this is the next
firmware change planned. A second wireless channel
(\mbox{nRF24L01+}) was scoped as a redundancy/comparison feature but
deferred. On the mechanical side, a force-sensitive gripper would
allow grasp-quality feedback, which is currently absent.

\section{Conclusion}
\label{sec:conclusion}

A wireless 4-DOF robotic arm with gripper has been designed,
implemented and is being evaluated against the EE6003 brief. The
two-unit embedded architecture, closed-loop joint control,
custom-framed Bluetooth protocol and Sense--Think--Act autonomous
mode together address every core objective, while the explicit
advanced features listed above provide the evidence base for marks
above the upper-2:1 ceiling. The system has been documented from
schematic through to firmware and test methodology in this report,
and a demonstration video accompanies the submission to evidence
prototype operation.
